\section*{Capítulo 6: Agujeros negros: los campeones de la curvatura}
\begin{description}[style=multiline,leftmargin=2em,labelwidth=1.5em]

% DOCX paragraph 2057
\item[1.]\hypertarget{note:06-01}{}Durante algún tiempo, el nombre de «\emph{agujero negro}» tuvo un rival: «\emph{estrella congelada}». Este último se acuñó para subrayar que, cuando una estrella colapsa y da origen a un agujero negro, para un observador externo parece «congelarse» en el tiempo. Sin embargo, la rivalidad no duró mucho, y la comunidad científica adoptó enseguida la expresión que todos conocemos hoy, tan eficaz como evocadora.

% DOCX paragraph 2058
\item[2.]\hypertarget{note:06-02}{}En cambio, si el agujero negro tuviera una masa del orden de unas pocas masas solares, las fuerzas de marea\index{fuerzas de marea}, es decir, las fuerzas gravitatorias generadas por la rápida variación del campo gravitatorio entre un punto y otro, serían enormemente distintas incluso entre tu cabeza y tus pies. Por eso, te estirarías como una goma elástica y acabarías convertido en un «espagueti». Francamente, es una experiencia que no recomiendo, ni siquiera como experimento mental~\ldots{}

% DOCX paragraph 2059
\item[3.]\hypertarget{note:06-03}{}J. Michell, en \emph{Philosophical Transactions of the Royal Society of London}, 74, 35 (1784).

% DOCX paragraph 2060
\item[4.]\hypertarget{note:06-04}{}R. P. Kerr, en \emph{Physical Review Letters}, 11, 237 (1963).

% DOCX paragraph 2061
\item[5.]\hypertarget{note:06-05}{}R. Ruffini y J. A. Wheeler, en \emph{Physics Today} 24, 30 (1971).

% DOCX paragraph 2062
\item[6.]\hypertarget{note:06-06}{}S. W. Hawking, en \emph{Communications of Mathematical Physics}, 43, 199 (1975); S. W. Hawking, en \emph{Nature} 248, 3031 (1974).

% DOCX paragraph 2063
\item[7.]\hypertarget{note:06-07}{}Como hemos dicho, los agujeros negros supermasivos tienen campos gravitatorios que varían muy gradualmente a escala humana, es decir, para «objetos» de dimensiones del orden de un par de metros. Por eso he elegido uno para el ejemplo.

% DOCX paragraph 2064
\item[8.]\hypertarget{note:06-08}{}Las dos características no están necesariamente ligadas: hay soluciones de agujeros negros que tienen un horizonte de sucesos, pero carecen de singularidad física\index{singularidad!física} en el centro. Por eso se llaman \emph{agujeros negros regulares}.

% DOCX paragraph 2065
\item[9.]\hypertarget{note:06-09}{}P. O. Mazur y E. Mottola, en \emph{Proceedings of the National Academy of Sciences, USA}, 101, 9545 (2004).

% DOCX paragraph 2066
\item[10.]\hypertarget{note:06-10}{}C. Misner y J. A. Wheeler, en \emph{Annals of Physics}, 2, 525 (1957).

\end{description}
